% 0722(본책 128쪽) — 제기차기 대회 참가자 20명의 1차(가로 8~16)·2차(세로 8~16) 성공 개수 산점도(격자 1개 · 눈금 2개 간격). 스캔 산점도를 TikZ 로 다시 그림(점 20개 원문 위치).
\begin{tikzpicture}[line width=0.5pt]
  \draw[gray!55, line width=0.25pt] (0,0) grid[xstep=0.27cm, ystep=0.27cm] (2.7,2.7);
  \draw[->, >=stealth, line width=0.7pt] (0,0) -- (2.85,0);
  \draw[->, >=stealth, line width=0.7pt] (0,0) -- (0,2.85);
  \draw[line width=0.5pt] (0.085,-0.06) -- (0.145,0.06) (0.155,-0.06) -- (0.215,0.06);
  \draw[line width=0.5pt] (-0.06,0.085) -- (0.06,0.145) (-0.06,0.155) -- (0.06,0.215);
  \node[font=\scriptsize, below left=0pt, inner sep=1.5pt] at (0,0) {$0$};
  \node[font=\scriptsize, below=1pt, inner sep=1pt] at (0.27,0) {$8$};
  \node[font=\scriptsize, below=1pt, inner sep=1pt] at (0.81,0) {$10$};
  \node[font=\scriptsize, below=1pt, inner sep=1pt] at (1.35,0) {$12$};
  \node[font=\scriptsize, below=1pt, inner sep=1pt] at (1.89,0) {$14$};
  \node[font=\scriptsize, below=1pt, inner sep=1pt] at (2.43,0) {$16$};
  \node[font=\scriptsize, left=1pt, inner sep=1pt] at (0,0.27) {$8$};
  \node[font=\scriptsize, left=1pt, inner sep=1pt] at (0,0.81) {$10$};
  \node[font=\scriptsize, left=1pt, inner sep=1pt] at (0,1.35) {$12$};
  \node[font=\scriptsize, left=1pt, inner sep=1pt] at (0,1.89) {$14$};
  \node[font=\scriptsize, left=1pt, inner sep=1pt] at (0,2.43) {$16$};
  \node[font=\scriptsize, anchor=north east, inner sep=1pt] at (2.85,-0.38) {1차(개)};
  \node[font=\scriptsize, align=center, anchor=north east, inner sep=1pt] at (-0.42,2.85) {2\\차\\(개)};
  \fill (0.27,0.54) circle (1.5pt);
  \fill (0.54,1.62) circle (1.5pt);
  \fill (0.54,1.08) circle (1.5pt);
  \fill (0.54,0.27) circle (1.5pt);
  \fill (0.81,1.08) circle (1.5pt);
  \fill (0.81,0.27) circle (1.5pt);
  \fill (1.08,1.62) circle (1.5pt);
  \fill (1.08,1.08) circle (1.5pt);
  \fill (1.08,0.81) circle (1.5pt);
  \fill (1.35,1.35) circle (1.5pt);
  \fill (1.35,0.54) circle (1.5pt);
  \fill (1.62,2.16) circle (1.5pt);
  \fill (1.62,1.62) circle (1.5pt);
  \fill (1.62,1.08) circle (1.5pt);
  \fill (1.62,0.81) circle (1.5pt);
  \fill (1.89,1.62) circle (1.5pt);
  \fill (2.16,1.89) circle (1.5pt);
  \fill (2.16,1.08) circle (1.5pt);
  \fill (2.43,2.43) circle (1.5pt);
  \fill (2.43,2.16) circle (1.5pt);
\end{tikzpicture}
